\documentclass[conference]{IEEEtran}
\usepackage{cite}
\usepackage{amsmath,amssymb,amsfonts}
\usepackage{algorithm}
\usepackage{algorithmic}
\usepackage{graphicx}
\usepackage{textcomp}
\usepackage{xcolor}
\usepackage{booktabs}
\usepackage{hyperref}

\begin{document}

\title{Proactive Comfort Agent: Full Self-Driving Home Orchestration\\
\large Technical Specification Document (Phase 2 -- Deliverable B)}

\author{
\IEEEauthorblockN{Ha Yun Kim}
\IEEEauthorblockA{\textit{Dept. of Information Systems} \\
\textit{Hanyang University}\\
Seoul, Korea \\
hykim1215@hanyang.ac.kr}
\and
\IEEEauthorblockN{Soo Young Song}
\IEEEauthorblockA{\textit{Dept. of Information Systems} \\
\textit{Hanyang University}\\
Seoul, Korea \\
sys1109@hanyang.ac.kr}
\and
\IEEEauthorblockN{Ruitan Sun}
\IEEEauthorblockA{\textit{Dept. of Information Systems} \\
\textit{Hanyang University}\\
Seoul, Korea \\
sunruitan2024@126.com}
\and
\IEEEauthorblockN{Tang Bin}
\IEEEauthorblockA{\textit{Dept. of Information Systems} \\
\textit{Hanyang University}\\
Seoul, Korea \\
tbrobo99@outlook.com}
}

\maketitle

\begin{abstract}
Modern smart homes present usability and sustainability barriers when users must manage disconnected appliance apps, monitor indoor environmental hazards, and coordinate high-wattage operations. We propose \textbf{FSD-Home (Full Self-Driving Home Orchestration)}, a proactive orchestration system designed for inclusive, zero-labor, and energy-conscious living. By evaluating real-time IoT telemetry streams (ambient temperature, humidity decay, appliance power draw, and compressor duty cycles) alongside household schedules, FSD-Home predicts indoor environmental and energy risks---such as post-wash mold buildup, kitchen thermal spikes, and peak electricity surges---before they escalate. To ensure safety, sustainability, and user autonomy, the system introduces a Tesla-inspired Risk-Graded Framework: Level 1 eco-efficiency and environmental adjustments run automatically with plain-language updates, while Level 2 high-load or schedule-altering operations require explicit, per-action user approval via a touch-first, accessible interface overlay backed by an append-only audit trail.
\end{abstract}

\begin{IEEEkeywords}
Smart Home Orchestration, FSD-Home, Full Self-Driving Home, Risk-Graded Autonomy, Energy Load-Balancing, Environmental Sustainability, Level 1 and Level 2 Control, Functional Safety, LG ThinQ Ecosystem.
\end{IEEEkeywords}

\section{Introduction}

\subsection{Motivation}
Smart home appliances are ubiquitous today, but having more features does not automatically ensure a seamless or energy-efficient experience for every member of a household \cite{msit2026}. Complex menus, fragmented mobile applications, and lack of cross-device coordination create unnecessary cognitive burden \cite{moneytoday2026}, with over 67\% of surveyed consumers reporting friction with hidden menus and unclear controls \cite{maeil2026}. At the same time, routine household oversights cause significant safety, indoor air quality, and energy-waste issues: 35.3\% of residential fires start from unattended cooking \cite{seoulfire2020, nfa2025}; wet laundry left in closed washers quickly leads to mold and bacterial growth requiring energy-intensive sanitizing wash cycles \cite{epa2022, who2009}; and uncoordinated high-wattage appliance use drives peak-hour electricity surges and wasted HVAC cooling.

\subsection{Problem Statement}
Current smart home apps like LG ThinQ work in reactive silos \cite{lgspec2026}. They wait for manual button taps or send passive alerts (such as ``Cycle Finished''), leaving household members to diagnose environmental or energy issues themselves \cite{pradhan2020}. More importantly, appliances do not share thermodynamic or power-load awareness across rooms. When someone bakes in a hot oven while running a dryer, kitchen temperatures surge, forcing the refrigerator compressor and air conditioner to run at maximum duty cycles and draw peak electricity without coordinated pre-cooling, ventilation, or load interruption \cite{ashrae2023}. FSD-Home solves this by combining continuous sensor tracking with automated cross-appliance environmental and energy orchestration.

\subsection{Related Software}
Current commercial tools offer strong single-device features: LG AI DD\textsuperscript{TM} senses fabric softness via motor torque; LG ezDispense\textsuperscript{TM} automates detergent dosing; and LG ThinQ ON / ThinQ Claw provides conversational voice hubs \cite{lgwhitepaper2024, lgspec2026}. Samsung's Bespoke AI sends lifestyle pattern notifications \cite{samsung2026a, samsung2026b}. However, none of them coordinate thermal load, moisture prevention, and electricity demand across multiple devices \cite{yuan2024}. FSD-Home introduces three main contributions: (1) explicit predictive risk scores for indoor mold hygiene (MRI) and kitchen thermal/power overload (TOI), (2) a Tesla-inspired Level 1 / Level 2 autonomy model with per-action confirmation, energy-saving pre-cooling/load interruption, and safe standby fallback, and (3) an inclusive, touch-first overlay designed around WCAG 2.2 contrast and touch targets \cite{wcag2023}.

\section{Requirements Overview}

\subsection{Functional Requirements (FR)}
\begin{itemize}
    \item \textbf{FR-1 (Real-Time Environmental \& Power Telemetry Perception):} Continuously sample indoor temperature, humidity, drum moisture decay, door contact, household presence (\texttt{user\_home}), compressor duty cycle, and wattage at 1~Hz.
    \item \textbf{FR-2 (Predictive Environmental Hazard \& Energy Modeling):} Deterministically compute the Mold Risk Index (MRI) and Thermal Overload Index (TOI) to trigger interventions at a threshold of $50.0$.
    \item \textbf{FR-3 (Level 1 \& Level 2 Action Framework):} Enforce a strict 7-action whitelist separating automatic eco/comfort adjustments (Level 1) from per-action approval-gated interventions (Level 2).
    \item \textbf{FR-4 (Inclusive UI/UX \& Per-Action Consent Overlay):} Provide WCAG 2.2 compliant touch controls, separate approval prompts per Level 2 action, supplementary TTS voice prompts, manual quiet-period status review, and truthful adapter-grounded result messages.
    \item \textbf{FR-5 (Fail-Safe Default \& Pre-Execution Safety Policy):} Enforce a 300-second approval expiration defaulting to safe standby, pre-execution state/risk revalidation, cooktop presence and auto-off safeguards, and dryer restart reminders.
    \item \textbf{FR-6 (Structured Append-Only Audit Trail):} Maintain an immutable SQLite \texttt{audit\_log} linking risk triggers, per-action approvals (\texttt{approval\_id}), and actual device execution outcomes.
\end{itemize}

\subsection{Non-Functional Requirements (NFR)}
\begin{itemize}
    \item \textbf{NFR-1 (Latency \& Architectural Decoupling):} Complete deterministic local risk evaluation within $1.2\text{ ms}$ per 1~Hz tick ($30\text{ ms}$ tick budget) while isolating cloud LLM reasoning ($2.0\text{--}3.0\text{ s}$) behind a 180-second cooldown and rule-based fallback.
    \item \textbf{NFR-2 (Simulation Containment \& Verification):} Validate system behavior locally without physical appliance hardware using deterministic JSON/CSV telemetry streams, automated unit tests, and latency benchmarks.
\end{itemize}

\section{Development Environment}

\subsection{Choice of Software Development Platform}
The FSD-Home backend system is developed using \textbf{Python 3.10+} (targeting standard CPython 3.11 execution) without heavy external framework dependencies for core processing.
\begin{itemize}
    \item \textbf{Backend Runtime:} Built with Python standard library modules (\texttt{http.server}, \texttt{threading}, \texttt{sqlite3}) to maintain lightweight execution, high portability, and fast cold-start performance.
    \item \textbf{Frontend UI:} Developed using vanilla HTML5, CSS3, and ES6 JavaScript (no-build step, zero external dependencies). It interacts with the backend via HTTP REST endpoints (\texttt{GET /api/proactive}, \texttt{GET /api/approvals}, \texttt{POST /api/approvals/<id>/confirm}).
    \item \textbf{Storage and Persistence:} SQLite 3.43 utilizing Write-Ahead Logging (WAL) mode for transactional isolation, audit decision logs, and user preference state tracking.
    \item \textbf{Target Operating Environment:} Tested locally on Windows 11 (AMD Ryzen 7 5800H, 16\,GB RAM) without requiring discrete GPU acceleration. Designed for cross-platform portability across Linux (Ubuntu 22.04 LTS) and ARM-based single-board devices (e.g., Raspberry Pi 4, OpenWrt routers).
    \item \textbf{Cloud Hybrid Architecture:} While real-time risk checks run deterministically on local infrastructure in $1.2\text{ ms}$, LLM reasoning calls out asynchronously to OpenAI-compatible cloud endpoints (Zhipu GLM-4-Flash / vLLM API).
    \item \textbf{Version Control:} Managed via Git and hosted on GitHub (\url{https://github.com/tbrobo99/proactive-senior-care-comfort-agent}).
\end{itemize}

\subsection{Software and Algorithms in Use}
The system decouples real-time physical risk perception from LLM decision reasoning:
\begin{itemize}
    \item \textbf{Deterministic Risk Engines (\texttt{app/risk.py}):} Evaluates Mold Risk Index (MRI) and Thermal Overload Index (TOI) using closed-form thermodynamic models and 30-minute least-squares linear trend forecasting.
    \item \textbf{OpenAI-Compatible LLM Inference:} Integrates with Zhipu AI \texttt{glm-4-flash} (or local vLLM endpoints) for intent resolution and action candidate selection.
    \item \textbf{Deterministic Safety Gating (\texttt{app/policy.py}):} Restricts LLM output via a strict 7-action candidate whitelist divided into Level 1 (Autonomous) and Level 2 (Approval Required) tiers. Unknown or malformed action names are rejected and logged rather than executed.
\end{itemize}

\subsection{Cost Estimation}
\begin{table}[htbp]
\caption{Cost Estimation}
\begin{center}
\begin{tabular}{llr}
\toprule
\textbf{Item} & \textbf{Specification} & \textbf{Cost} \\
\midrule
Development HW & Existing Ryzen 7 PC & \textwon 0 \\
Core Software & Python 3.11, SQLite (Open-source) & \textwon 0 \\
Cloud API / LLM & Zhipu GLM-4-Flash / vLLM API & \textwon 0 (Free version) \\
Deployment & Local / Edge Simulation & \textwon 0 \\
\midrule
\textbf{Total Estimated Cost} & & \textbf{\textwon 0} \\
\bottomrule
\end{tabular}
\end{center}
\end{table}

\subsection{Task Distribution}
\begin{table}[htbp]
\caption{Role Assignments}
\begin{center}
\begin{tabular}{llp{4.3cm}}
\toprule
\textbf{Role} & \textbf{Assigned Member} & \textbf{Task Description \& Responsibilities} \\
\midrule
User & Ruitan Sun & Identify user needs, provide feedback, and help ensure the service meets user expectations. \\
Customer & Soo Young Song & Functional safety boundaries, Level 2 approval threshold definitions, audit log schema. \\
Software Developer & Tang Bin & General backend part: Agentic loop (\texttt{proactive.py}), risk engine (\texttt{risk.py}), telemetry generator (\texttt{sensors.py}). \\
Development Manager & Ha Yun Kim & System architecture design, LaTeX document formatting, UI design and accessibility validation. \\
\bottomrule
\end{tabular}
\end{center}
\end{table}

\section{Specifications}

\subsection{System Architecture Overview}
Fig.~\ref{fig:arch} illustrates the decoupled FSD-Home pipeline, separating 1~Hz deterministic telemetry evaluation and safety whitelisting from asynchronous cloud LLM ranking, per-action Level 2 approval modals, and append-only audit logging.

\begin{figure}[htbp]
\centering
\fbox{\parbox{0.94\columnwidth}{\footnotesize\ttfamily
[1 Hz Telemetry Stream: app/sensors.py]\\
Temp | Humidity | Door | Power (W) | Fridge Duty | user\_home\\
\hspace*{3em}$\downarrow$ (1.2 ms deterministic check)\\
{[Deterministic Risk \& Energy Engine: app/risk.py]}\\
MRI >= 50.0 (Mold Risk) \quad | \quad TOI >= 50.0 (Thermal/Load)\\
\hspace*{3em}$\downarrow$\\
{[Safety Whitelist \& Policy Gate: app/policy.py]}\\
Reject Unknown Actions | Assign Level 1 vs. Level 2\\
\hspace*{1em}$\swarrow$ (Level 1: Eco/Comfort) \hspace*{2em} $\searrow$ (Level 2: High-Load/Heat)\\
{[LLM Ranker / Abstain]} \hspace*{3.5em} {[Per-Action Approval Queue]}\\
(2--3 s; 180 s cooldown) \hspace*{3.5em} (app/approval.py, 300s TTL)\\
\hspace*{1em}$\searrow$ \hspace*{10em} $\swarrow$ [Confirm] + Revalidate\\
{[Device Adapters \& Append-Only Audit Log: app/db.py]}\\
Write Intent $\rightarrow$ Execute Adapter $\rightarrow$ Log Outcome $\rightarrow$ Update UI
}}
\caption{FSD-Home decoupled sensing, safety gating, per-action Level 2 approval, and append-only audit architecture.}
\label{fig:arch}
\end{figure}

\subsection{FR-1: Real-Time Environmental \& Power Telemetry Perception}
The \texttt{SensorSimulator} class in \texttt{app/sensors.py} generates 1~Hz time-series IoT telemetry snapshots.
\begin{itemize}
    \item \textbf{Data Pipeline:} Every snapshot is appended to \texttt{data/sensors/readings.jsonl}, and the latest readings are rewritten to \texttt{data/sensors/latest.csv}.
    \item \textbf{Simulated Metrics:} Room temperatures ($^\circ$C), relative humidity (\%), drum moisture decay, appliance door contact, household presence (\texttt{user\_home}), refrigerator compressor duty cycle, and active power draw (W).
\end{itemize}

\subsection{FR-2: Predictive Risk \& Energy Modeling}
The risk assessment module (\texttt{app/risk.py}) computes MRI and TOI deterministically within $1.2\text{ ms}$ per tick to protect indoor air quality and prevent wasteful peak-load energy spikes:
\begin{equation}
\text{MRI} = \text{clamp}\left[(H - 70.0) \times 3.2 + (S \times 6.0) + \max(0, T - 24.0) \times 1.5\right]
\end{equation}
\begin{equation}
\text{TOI} = \text{clamp}\left[\frac{W}{25.0} + \max(0, T_k - 27.0) \times 6.0 + \max(0, D - 0.6) \times 60.0\right]
\end{equation}
\begin{itemize}
    \item \textbf{Environmental \& Hygiene Optimization (MRI):} Prevents mold proliferation in sealed appliances early using low-wattage ventilation or dehumidification, avoiding harsh chemical cleaners and energy-heavy hot-water tub-clean cycles \cite{epa2022, who2009}.
    \item \textbf{Thermodynamic \& Energy Load Balancing (TOI):} Tracks concurrent heat wattage ($W$), kitchen temperature ($T_k$), and refrigerator compressor duty cycle ($D > 0.6$, signaling excess thermal strain) alongside scheduled cooking events \cite{ashrae2023}.
    \item \textbf{Trigger \& Tier Separation:} When either index reaches or exceeds its configured threshold ($50.0$), an intervention event is dispatched. The risk threshold determines \emph{whether} an intervention is considered, whereas the predefined action tier determines \emph{whether} the action runs automatically or requires user approval.
\end{itemize}

\subsection{FR-3: Level 1 \& Level 2 Action Framework}
The local rule engine (\texttt{app/policy.py}) strictly enforces a 7-action whitelist (Table~\ref{tab:whitelist}). Unknown actions are blocked and recorded in \texttt{audit\_log}.
\begin{itemize}
    \item \textbf{Level 1 (Autonomous Eco \& Comfort Actions):} Low-power, reversible actions execute immediately when selected---including proactive kitchen pre-cooling (\texttt{precool\_kitchen} prior to scheduled cooking so the AC compressor runs smoothly at high efficiency rather than surging against peak oven heat) and targeted ventilation/dehumidification.
    \item \textbf{Level 2 (Per-Action Confirmation Required):} High-load or schedule-interrupting operations---such as stopping the dryer (\texttt{stop\_dryer}) during thermal/power overload so high-wattage heating elements do not overlap with kitchen cooking---generate individual pending confirmation states in \texttt{app/approval.py}. In the current proactive flow, Level 2 interventions are triggered by thermal/power overload risk (TOI), whereas mold risk (MRI) produces Level 1 actions.
\end{itemize}

\begin{table}[htbp]
\caption{Proposed Whitelist}
\label{tab:whitelist}
\begin{center}
\begin{tabular}{lll}
\toprule
\textbf{Tier} & \textbf{Action} & \textbf{Device} \\
\midrule
L1 & \texttt{ventilate\_washer} & Washer, drum vent \\
L1 & \texttt{dehumidifier\_on} & Balcony dehumidifier \\
L1 & \texttt{purifier\_on} & Living-room purifier \\
L1 & \texttt{ventilate\_kitchen} & Range hood, fan high \\
L1 & \texttt{precool\_kitchen} & Kitchen AC to $24^\circ\text{C}$ \\
L2 & \texttt{stop\_dryer} & Dryer off; restart by hand \\
L2 & \texttt{preheat\_cooktop} & Cooktop on early \\
\bottomrule
\end{tabular}
\end{center}
\end{table}

\subsection{FR-4: Inclusive UI/UX and Interface Specification}
\begin{itemize}
    \item \textbf{Per-Action Level 2 Approval Modal:} Rendered dynamically via REST polling (\texttt{GET /api/approvals}). Each Level 2 action is presented with its own separate plain-language question and unique \texttt{approval\_id} so approving one action cannot authorize another.
    \item \textbf{Visual \& Touch Targets:} Text contrast exceeds $4.5:1$, and primary touch targets are $\ge 48 \times 48\text{ px}$ with clearly labeled \texttt{[Confirm]} and \texttt{[Decline]} buttons.
    \item \textbf{Real-Time Visual Countdown:} Displays a visible 300-second countdown timer. Upon timeout, the request transitions to \texttt{expired} without altering appliance state.
    \item \textbf{Touch-First Multimodal Support:} Touch buttons serve as the primary consent method. Spoken prompts (via Web Speech API TTS) and visible text act as supplementary aids; silence, ambiguous speech, or prompt playback alone never count as consent.
    \item \textbf{Truthful Outcome Reporting \& Quiet-Period Status View:} Reports actual adapter outcomes (\texttt{completed}, \texttt{canceled}, or \texttt{failed}) with a concise reason, keeping technical telemetry in the audit log. During the 30-minute quiet period after a decline or timeout, users can still open the interface to manually check live risk status.
\end{itemize}

\subsection{FR-5: Fail-Safe Default Policy}
\begin{itemize}
    \item \textbf{300-Second Expiry \& Safe Standby:} Level 2 confirmation requests expire after 300 seconds (\texttt{APPROVAL\_TTL\_SECONDS}). Decline, expiry, late responses, or system restarts never execute pending Level 2 actions and initiate a 1,800-second (30-minute) quiet period for that risk kind.
    \item \textbf{Pre-Execution Revalidation \& Cooktop Safeguards:} Upon confirmation within 300 seconds, the engine re-verifies approval validity, active risk conditions, device state, and safety interlocks immediately before calling the device adapter. Because \texttt{preheat\_cooktop} prepares a scheduled meal rather than relieving the overload, it is revalidated against the meal instead of the risk score: execution requires verified household presence (\texttt{user\_home}), a cooking session still scheduled within 30 minutes, and the cooktop still off. After execution, an automatic shut-off interlock switches the cooktop off if cooking is not detected within 10 minutes of the scheduled start (preventing wasted energy and fire hazards), and in any case after a 90-minute maximum continuous run time.
    \item \textbf{Follow-Up Notification:} When \texttt{stop\_dryer} interrupts a cycle to relieve thermal and electrical overload, the system notifies the user once the risk clears so they can restart the dryer manually.
\end{itemize}

\subsection{FR-6: Structured Append-Only Audit Trail}
To ensure end-to-end traceability across risk evaluation, user consent, and hardware execution (\texttt{app/db.py}), the system maintains an append-only SQLite \texttt{audit\_log} table (protected against \texttt{UPDATE} and \texttt{DELETE} via database triggers):
\begin{itemize}
    \item \textbf{Linked Event Records:} Captures timezone-aware ISO 8601 timestamps, \texttt{event\_type} (e.g., \texttt{risk.triggered}, \texttt{approval.proposed}, \texttt{approval.confirmed}, \texttt{action.executed}, \texttt{action.failed}, \texttt{action.blocked}), \texttt{risk\_kind}, \texttt{risk\_score}, \texttt{action\_kind}, \texttt{tier}, \texttt{approval\_id}, \texttt{approval\_status}, \texttt{execution\_status}, \texttt{decision\_source}, and \texttt{channel}.
    \item \textbf{Write-Before-Act Integrity:} Action intent is recorded prior to invoking the device adapter, and actual adapter outcomes are logged immediately afterward. If audit persistence fails prior to a Level 2 execution, the action is aborted.
\end{itemize}

\subsection{Control Flow and Algorithmic Execution}
Algorithm~\ref{alg:control} defines the 1~Hz orchestration loop, per-action Level 2 approval queuing, and pre-execution revalidation.

\begin{algorithm}[htbp]
\caption{1~Hz Proactive Orchestration and Per-Action Approval}
\label{alg:control}
\begin{algorithmic}[1]
\STATE \textbf{function} \textsc{tick}() \COMMENT{1 Hz sampling; decisions may wait}
\STATE \quad $snap \leftarrow \text{sensors.snapshot}()$; $\text{sensors.record}(snap)$
\STATE \quad \textbf{for} $r \in [r \in \text{risk.evaluate}(snap, hist) \mid r.\text{score} \ge 50.0]$ \textbf{do}
\STATE \quad \quad \textbf{if} $\text{cooling}(r.\text{kind}) \lor \text{pending}(r.\text{kind})$ \textbf{then continue}
\STATE \quad \quad $\textsc{intervene}(r, snap)$ \COMMENT{may call the model}
\STATE \quad \textbf{for} $ap \in \text{approvals.expire\_due}()$ \textbf{do} \COMMENT{300s deadline check}
\STATE \quad \quad $\text{quiet}(ap.\text{risk\_kind}, 30\text{ min})$; $\text{report}(ap)$
\STATE
\STATE \textbf{function} \textsc{intervene}($r, snap$)
\STATE \quad $cands \leftarrow \text{policy.candidates}(r, snap)$ \COMMENT{whitelist}
\STATE \quad $chosen \leftarrow \text{decide}(r, cands)$ \COMMENT{LLM ranks L1 or abstains}
\STATE \quad \textbf{for} $c \in cands$ \textbf{do}
\STATE \quad \quad \textbf{if} $\text{tier\_of}(c) == \text{TIER2}$ \textbf{then}
\STATE \quad \quad \quad $\text{approval.propose}(c, \text{ttl}=300)$ \COMMENT{per-action ID}
\STATE \quad \quad \textbf{else if} $c \in chosen \land \text{revalidate}(c, snap)$ \textbf{then}
\STATE \quad \quad \quad $\text{run\_and\_log}(c)$
\STATE
\STATE \textbf{function} \textsc{on\_confirm}($approval\_id, channel$)
\STATE \quad $ap \leftarrow \text{approval.confirm}(approval\_id, channel)$
\STATE \quad \textbf{if} $ap \land \text{revalidate}(ap.\text{action}, \text{sensors.snapshot}())$ \textbf{then}
